\chapter{Strumenti di front-end}

Abbiamo iniziato dalle tecnologie fondamentali (HTML, CSS, JavaScript),
costruito pagine statiche e poi siamo passati alla programmazione lato server.
Le pagine che generavamo dinamicamente restavano però semplici: un solo foglio
di stile, pochissimo JavaScript lato client, nessuna libreria esterna.

I front-end moderni sono molto più complessi:

\begin{itemize}
  \item il codice JavaScript lato client è sempre più articolato;
  \item si usano molte librerie e moduli;
  \item si transpila e si applica il \textit{downleveling};
  \item si ottimizzano i file serviti;
  \item si usano spesso pre- e post-processori per i fogli di stile.
\end{itemize}

Servono strumenti per assistere lo sviluppatore e automatizzare questi
passaggi.

\section{Pre- e post-processori CSS}

Il CSS di per sé può essere piacevole, ma i fogli di stile diventano presto
grandi, complessi e difficili da mantenere. Gli strumenti di pre- e
post-elaborazione permettono di creare in modo efficiente fogli di stile
robusti, scalabili e ottimizzati.

\begin{figure}[H]
  \centering
  \begin{tikzpicture}[font=\Lato\small,
      n/.style={wtnode, minimum width=2.9cm, minimum height=1.0cm},
      f/.style={wtnodeplain, minimum width=2.2cm, minimum height=0.9cm},
      a/.style={-{Stealth[length=5pt,width=4pt]}, line width=0.85pt,
                draw=neutral!45!black}]

    \node[f] (src) at (0,0) {file sorgente};
    \node[n, draw=orange-gradient-6, fill=orange-gradient-1!40!white]
      (pre) at (3.4,0) {\textbf{Pre-processore}};
    \node[f] (c1) at (6.5,0) {CSS};
    \node[n, draw=azure-gradient-6, fill=azure-gradient-1!45!white]
      (post) at (9.6,0) {\textbf{Post-processore}};
    \node[f] (c2) at (12.6,0) {CSS};

    \draw[a] (src) -- (pre);
    \draw[a] (pre) -- (c1);
    \draw[a] (c1) -- (post);
    \draw[a] (c2.south) -- ++(0,-0.9) node[below, font=\Lato\scriptsize,
                                          text=neutral!30!black] {browser};
    \draw[a] (post) -- (c2);
  \end{tikzpicture}
  \caption{La catena di elaborazione di un foglio di stile moderno.}
  \label{fig:css-pipeline}
\end{figure}

\subsection{Pre-processori}

I \term{pre-processori CSS} sono \guillemotleft compilatori\guillemotright\ speciali che generano
codice CSS a partire da file scritti in una sintassi propria. Tale sintassi si
può vedere come un'\textbf{estensione} del CSS di base, che introduce
funzionalità nuove (per esempio le variabili) per scrivere fogli di stile
manutenibili in modo più efficiente.

I più diffusi sono \term{Sass}, \term{LESS} e \term{Stylus}.

\subsection{Post-processori}

I \term{post-processori} elaborano e trasformano automaticamente fogli di stile
già esistenti. Si usano tipicamente per:

\begin{itemize}
  \item \textbf{downleveling}, cioè riscrivere costrutti recenti in forme
        comprese da browser più vecchi:

\begin{lstlisting}[style=css]
/* prima */                  /* dopo */
body {                       body {
  color: oklch(61% 0.2 29);    color: rgb(225, 65, 52);
}                            }
\end{lstlisting}

  \item \textbf{autoprefissatura}, cioè aggiungere i prefissi specifici dei
        browser:
\end{itemize}

\begin{lstlisting}[style=css]
/* prima */                  /* dopo */
.example {                   .example {
  transition: all .5s;         -webkit-transition: all .5s;
}                              -o-transition: all .5s;
                               transition: all .5s;
                             }
\end{lstlisting}

\begin{itemize}
  \item \textbf{ottimizzazioni}, come la minificazione o l'eliminazione delle
        parti inutilizzate;
  \item \textbf{imposizione di convenzioni} e rilevamento di errori.
\end{itemize}

Il post-processore più diffuso è \term{PostCSS}.

\section{Sass}

\term{Sass} si autodefinisce \guillemotleft CSS con i superpoteri\guillemotright\ ed è attualmente il
pre-processore CSS più diffuso. Esiste dal 2006. Gli altri pre-processori
offrono grosso modo le stesse funzionalità con una sintassi leggermente diversa.

\begin{lstlisting}[style=shell]
npm install -g sass
\end{lstlisting}

\subsection{Due sintassi}

\begin{center}
\renewcommand{\arraystretch}{1.4}
\begin{tabular}{@{}>{\raggedright\arraybackslash}p{6.3cm}@{\hspace{1.0cm}}>{\raggedright\arraybackslash}p{6.3cm}@{}}
  \textbf{Sintassi SCSS} & \textbf{Sintassi indentata} \\
  Estensione \code{.scss}. & Estensione \code{.sass}. \\
  È un \textbf{superset} del CSS: qualunque CSS valido è SCSS valido.
  & Basata sull'indentazione, come Python. \\
  È la sintassi più diffusa. & Era la sintassi originale. \\
\end{tabular}
\end{center}

\subsection{Variabili}

Le variabili permettono di riusare valori in tutto il foglio di stile. Al momento
della compilazione vengono sostituite dal loro valore.

\begin{lstlisting}[style=css, name=file.scss]
$font-stack: Helvetica, sans-serif;
$primary-color: #333;

body {
  font: $font-stack;
  color: $primary-color;
}
\end{lstlisting}

\begin{lstlisting}[style=css, name=file.css]
body {
  font: Helvetica, sans-serif;
  color: #333;
}
\end{lstlisting}

\subsection{Annidamento}

In HTML esiste una chiara gerarchia di elementi annidati; nel CSS puro non
possiamo avere regole annidate. Sass lo permette, e rende il codice più
leggibile.

\begin{lstlisting}[style=css, name=nesting.scss]
nav {
  ul {
    margin: 0;
    padding: 0;
    list-style: none;
  }

  li { display: inline-block; }

  a {
    display: block;
    padding: 6px 12px;
    text-decoration: none;
  }
}
\end{lstlisting}

\begin{lstlisting}[style=css, name=nesting.css]
nav ul { margin: 0; padding: 0; list-style: none; }
nav li { display: inline-block; }
nav a  { display: block; padding: 6px 12px; text-decoration: none; }
\end{lstlisting}

\warning{Annidare troppo è una cattiva pratica}{
  Ogni livello di annidamento aggiunge un elemento al selettore generato, e
  quindi aumenta la \textbf{specificità}. Selettori troppo specifici sono
  difficili da sovrascrivere e rendono il foglio di stile fragile: torna qui
  tutto quello che abbiamo studiato sulla cascata nel capitolo 3. Come regola
  pratica, non superare i tre livelli.}

\subsection{Mixin}

I \term{mixin} permettono di riusare un gruppo di dichiarazioni. Si dichiarano
con la regola \code{@mixin} e si includono con \code{@include}.

\begin{lstlisting}[style=css, name=mixins.scss]
@mixin clear-list {
  list-style-type: none;
  padding: 0;
  margin: 0;
  li {
    display: inline-block;
  }
}

nav ul {
  background-color: aliceblue;
  @include clear-list;
}
\end{lstlisting}

I mixin possono anche avere parametri, con valori predefiniti:

\begin{lstlisting}[style=css, name=theme.scss]
@mixin theme($color: DarkBlue) {
  background: $color;
  color: #fff;
}

.info   { @include theme; }
.danger { @include theme($color: DarkRed); }
\end{lstlisting}

\subsection{Moduli e parziali}

\begin{lstlisting}[style=css, name=base.scss]
$font-stack: Helvetica, sans-serif;
$primary-color: #333;

body {
  font: $font-stack;
  color: $primary-color;
}
\end{lstlisting}

\begin{lstlisting}[style=css, name=example.scss]
@use 'base';

.inverse {
  background: base.$primary-color;
  color: white;
}
\end{lstlisting}

I file Sass il cui nome inizia con un trattino basso sono detti \term{parziali}:
sono pensati solo per essere inclusi da altri file, e Sass non genera per essi
un file CSS autonomo. Un file che raccoglie tutte le variabili è un buon esempio
di parziale.

\begin{lstlisting}[style=css, name=_variables.scss]
$primary-color: purple;
$accent-color: orange;
$background-color: beige;
\end{lstlisting}

\subsection{Liste, mappe e controllo di flusso}

\begin{lstlisting}[style=css]
$list: (120px, 240px, 480px, 800px);
$map: ("xs": 480px, "md": 768px, "lg": 1024px);
\end{lstlisting}

Sass offre regole per decidere \textbf{se} emettere degli stili, o per emetterli
più volte con piccole variazioni: si può pensarlo come un motore di template per
generare CSS. Sono disponibili \code{@if}/\code{@else}, \code{@each},
\code{@for} e \code{@while}.

\begin{lstlisting}[style=css, name=avatar.scss]
@use "sass:math";

@mixin avatar($size, $circle: false) {
  width: $size;
  height: $size;

  @if $circle {
    border-radius: math.div($size, 2);
  }
}

.square-av { @include avatar(100px, $circle: false); }
.circle-av { @include avatar(100px, $circle: true);  }
\end{lstlisting}

\begin{lstlisting}[style=css, name=hidden.scss]
$sizes: ("xs": 480px, "md": 768px, "lg": 1024px);

@each $sizeName, $sizeMaxWidth in $sizes {
  @media screen and (max-width: $sizeMaxWidth) {
    .hidden-#{$sizeName} {
      display: none;
    }
  }
}
\end{lstlisting}

Da tre righe di dichiarazione si generano tre media query complete: è
esattamente il tipo di ripetizione che rende ingestibili i fogli di stile
scritti a mano.

\section{Framework CSS}

Anche scrivendo il CSS del front-end ci si ritrova a risolvere sempre gli stessi
problemi: creare un sistema di impaginazione, definire una tipografia coerente,
dare stile a tabelle, liste e componenti, azzerare gli stili predefiniti dei
browser per ottenere un aspetto uniforme.

Esistono molti \term{framework CSS} che forniscono un insieme di stili di uso
generale, utilizzabili così come sono oppure come base da rifinire. Fra i più
diffusi: \term{Bootstrap} (il più popolare, originariamente di Twitter),
\term{Tailwind} (noto per l'approccio \textit{utility-first}),
\term{Materialize}, \term{Foundation}, \term{Bulma}, \term{PureCSS}.

\info{Non sono framework in senso stretto}{
  Nonostante il nome, non sono framework come quelli visti nel capitolo 13: non
  c'è inversione del controllo, non chiamano il nostro codice. Sono
  \textbf{insiemi di fogli di stile} da importare, con classi da applicare agli
  elementi.}

I framework sono uno o più file CSS: basta importarli nelle nostre pagine e
usare le classi che forniscono. Si possono scaricare e servire dal proprio
server, oppure si può ricorrere a una \term{CDN} (\textit{Content Delivery
Network}) che li ospita per noi --- comodo per prototipare, ma con qualche
implicazione sulla privacy, dato che ogni visitatore contatta un server di terze
parti.

\subsection{Personalizzare un framework}

Potremmo non volere che il nostro sito assomigli a tutti gli altri costruiti con
lo stesso framework, e gli stili predefiniti potrebbero non soddisfare
esigenze particolari. Ci sono due strade:

\begin{itemize}
  \item la maggior parte dei framework CSS è costruita con Sass o strumenti
        simili: si possono \textbf{cambiare le variabili} di build (colori,
        caratteri, dimensioni) e ricompilare il framework;
  \item si possono definire stili personalizzati che \textbf{sovrascrivono}
        alcuni stili del framework.
\end{itemize}

\begin{lstlisting}[style=css, name=custom.scss]
// sovrascrittura delle variabili di Bootstrap
$primary: #00647a;
$secondary: rgb(254, 203, 110);

@import "../node_modules/bootstrap/scss/bootstrap.scss";
\end{lstlisting}

Compilando con Sass si ottiene la propria versione personalizzata di Bootstrap.

\section{Bundling e minificazione}

\subsection{Bundling}

\dfn{Bundling}{
  Il \term{bundling} combina più file CSS o JavaScript in un unico file, così
  che il browser effettui \textbf{meno richieste HTTP} per caricare una
  pagina.}

Le richieste HTTP aggiuntive introducono latenza. Le strategie di caching
aiutano nei caricamenti successivi, ma includere molti fogli di stile e molti
script ha comunque un impatto sul \textbf{primo} caricamento: il bundling è
un'ottimizzazione ampiamente adottata.

\begin{lstlisting}[style=html, name=prima.html]
<link rel="stylesheet" href="a.css">
<link rel="stylesheet" href="b.css">
<!-- altri fogli di stile -->
<script type="javascript" src="a.js"></script>
<script type="javascript" src="b.js"></script>
<script type="javascript" src="c.js"></script>
<!-- altri script -->
\end{lstlisting}

\begin{lstlisting}[style=html, name=dopo.html]
<link rel="stylesheet" href="bundle.css">
<script type="javascript" src="bundle.js"></script>
\end{lstlisting}

\warning{L'ordine conta}{
  L'ordine in cui i singoli file vengono accodati nel bundle è significativo: per
  il CSS decide chi vince nella cascata a parità di specificità, per JavaScript
  decide quali definizioni sono disponibili al momento dell'esecuzione.}

\info{Una nota storica}{
  Prima che i moduli ES fossero disponibili nei browser non esisteva alcun modo
  di scrivere codice JavaScript modulare lato client. Strumenti come
  \term{Webpack}, \term{Rollup} e \term{Parcel} nacquero proprio per concatenare
  tutti i moduli sorgente in una forma che i browser potessero comprendere.}

\subsection{Tree shaking}

I \textit{bundler} eseguono spesso anche il \term{tree shaking}, cioè
l'eliminazione del codice morto.

Spesso ci serve solo una parte delle funzionalità dei moduli che importiamo, e i
moduli possono essere piuttosto grandi: includere un intero modulo quando ne
usiamo una sola funzione è inefficiente. Il tree shaking analizza importazioni ed
esportazioni per \textbf{potare} dal bundle finale tutto il codice inutilizzato,
con un impatto notevole sulla dimensione e quindi sulle prestazioni.

\begin{figure}[H]
  \centering
  \begin{tikzpicture}[font=\Lato\scriptsize,
      m/.style={rounded corners=3pt, draw=primary!60!black, line width=0.7pt,
                fill=primary!10!white, minimum width=2.6cm, minimum height=0.55cm},
      f/.style={rounded corners=2pt, draw=neutral!50!white, line width=0.6pt,
                fill=neutral!8!white, minimum width=1.15cm, minimum height=0.5cm},
      d/.style={f, draw=red-gradient-5, fill=red-gradient-1!30!white},
      a/.style={-{Stealth[length=4pt,width=3.5pt]}, line width=0.7pt,
                draw=neutral!45!black}]

    \node[f] (i) at (4.3,-1.6) {index.js};

    \node[m] (ma) at (1.4,-0.7) {Modulo A};
    \node[m] (mb) at (4.3,-0.7) {Modulo B};
    \node[m] (mc) at (7.2,-0.7) {Modulo C};

    \foreach \x/\lab/\st in {0.2/f1/f, 1.4/f2/f, 2.6/f3/d,
                             3.7/f4/d, 4.9/f5/f, 6.1/f6/f, 7.3/f7/f, 8.5/f8/f} {
      \node[\st] at (\x,0.4) {\lab};
    }

    \draw[a] (ma) -- (i);
    \draw[a] (mb) -- (i);
    \draw[a] (mc) -- (i);

    \node[font=\Lato\scriptsize, text=red-gradient-6!75!black, anchor=west]
      at (9.4,0.4) {codice inutilizzato:};
    \node[font=\Lato\scriptsize, text=red-gradient-6!75!black, anchor=west]
      at (9.4,0.0) {escluso dal bundle};
  \end{tikzpicture}
  \caption{Il tree shaking esclude dal bundle il codice mai raggiunto.}
  \label{fig:tree-shaking}
\end{figure}

\subsection{Minificazione}

Quando scriviamo codice, l'indentazione e i nomi significativi delle variabili
sono essenziali per leggibilità, comprensibilità e manutenibilità. Ai browser,
però, non serve codice leggibile: possiamo ottimizzare il trasferimento
rimuovendo gli spazi \guillemotleft superflui\guillemotright\ e accorciando il codice senza
alterarne il funzionamento. È la \term{minificazione}.

\begin{lstlisting}[style=js, name=originale.js (188 byte)]
let msg = "Hello";
let name = "Web Technologies";
let str = `${msg}, ${name}!`;
let b = true;
if (b === true)
  document.getElementById('msg').innerHTML = str;
\end{lstlisting}

\begin{lstlisting}[style=js, name=minificato.js (123 byte)]
let msg="Hello",name="Web Technologies",str=`${msg},${name}!`,b=!0;!0===b&&
(document.getElementById("msg").innerHTML=str);
\end{lstlisting}

Una compressione del 34,57\%, cioè 65 byte risparmiati. L'effetto è tanto
maggiore quanto più grandi sono i file sorgente:

\begin{center}
\renewcommand{\arraystretch}{1.4}
\begin{tabular}{@{}l@{\hspace{1.0cm}}r@{\hspace{1.0cm}}r@{\hspace{1.0cm}}r@{}}
  \textbf{Libreria} & \textbf{Originale} & \textbf{Minificata} &
  \textbf{Risparmio} \\
  React 18.2.0     & $\sim$85 kB  & $\sim$25 kB & $\sim$70\% \\
  Moment.js 2.29.4 & $\sim$151 kB & $\sim$75 kB & $\sim$50\% \\
\end{tabular}
\end{center}

\section{Vite}

\term{Vite} (si pronuncia \guillemotleft vit\guillemotright) è uno strumento che assiste lo
sviluppatore sia in fase di sviluppo sia in fase di rilascio. Ha due componenti
chiave:

\begin{itemize}
  \item un \textbf{server di sviluppo} con numerose funzionalità avanzate per
        snellire il lavoro;
  \item uno \textbf{strumento di build} per creare risorse statiche altamente
        ottimizzate, pronte per la produzione.
\end{itemize}

\begin{lstlisting}[style=shell]
npm create vite@latest
\end{lstlisting}

Il comando chiede il nome del progetto, il framework (Vanilla, Vue, React,
Preact, Lit, Svelte, Solid, Qwik) e la variante (JavaScript o TypeScript), poi
genera lo scheletro del progetto.

\begin{lstlisting}[style=json, name=package.json]
{
  "name": "hello-vite",
  "private": true,
  "version": "0.0.0",
  "type": "module",
  "scripts": {
    "dev": "vite",
    "build": "tsc && vite build",
    "preview": "vite preview"
  },
  "devDependencies": {
    "typescript": "^5.2.2",
    "vite": "^5.0.8"
  }
}
\end{lstlisting}

\begin{center}
\renewcommand{\arraystretch}{1.4}
\begin{tabular}{@{}l@{\hspace{1.2cm}}>{\raggedright\arraybackslash}p{8.8cm}@{}}
  \textbf{Comando} & \textbf{Effetto} \\
  \code{dev}     & avvia il server di sviluppo. Vite osserva le modifiche ai
                   file sorgente e ricarica automaticamente l'applicazione
                   (\textit{live reload}). \\
  \code{build}   & genera un bundle ottimizzato, pronto per la produzione. \\
  \code{preview} & è un semplice server statico che serve i file generati
                   durante la fase di build. \\
\end{tabular}
\end{center}

\subsection{Aggiungere Sass}

\begin{lstlisting}[style=shell]
npm install -D sass
\end{lstlisting}

Basta poi creare un foglio di stile Sass e importarlo nel file \code{main.ts}:
Vite rileva automaticamente che si tratta di un file Sass, lo compila e lo
include nel bundle dopo il file CSS originale.

\subsection{Che cosa fa la build}

\begin{console}[Terminale]
$ npm run build

> hello-vite@0.0.0 build
> tsc && vite build

vite v5.0.10 building for production...
 8 modules transformed.
dist/index.html                   0.46 kB | gzip: 0.29 kB
dist/assets/index-v5kLGx6_.css    1.21 kB | gzip: 0.63 kB
dist/assets/index-DiAK_ebR.js     3.05 kB | gzip: 1.62 kB
 built in 173ms
\end{console}

In una sola invocazione Vite ha:

\begin{itemize}
  \item analizzato quali risorse statiche sono effettivamente usate;
  \item compilato i file Sass in CSS e quelli TypeScript in JavaScript;
  \item riunito i file Sass e CSS in un unico file, e i file JavaScript in un
        altro;
  \item minificato i bundle;
  \item applicato il \term{fingerprinting} ai nomi dei file finali;
  \item spostato tutto nella directory \code{/dist}, aggiornando
        \code{index.html} con le importazioni corrette.
\end{itemize}

\subsection{Fingerprinting dei nomi}

I bundle prodotti hanno nomi curiosi: \code{index-v5kLGx6\_.css},
\code{index-DiAK\_ebR.js}. L'ultima parte del nome è un \textbf{hash del
contenuto}: costruendo una nuova versione dell'applicazione, se i file cambiano,
cambia anche il nome.

\info{Perché serve}{
  Supponiamo che i file si chiamassero semplicemente \code{index.css} e
  \code{index.js}. Un browser che visita la pagina un minuto prima del rilascio
  di una nuova versione continuerebbe a usare per un certo tempo le copie
  vecchie conservate nella propria cache, e vedrebbe un sito rotto. Cambiando
  il nome a ogni modifica, il problema scompare: la cache del vecchio file
  resta valida, ma nessuno lo richiede più.}

\section{Riferimenti}

\begin{itemize}
  \item \textbf{Client-side Tooling Overview} --- MDN web docs. \\
        \urlx{https://developer.mozilla.org/en-US/docs/Learn/Tools_and_testing/Understanding_client-side_tools/Overview}
  \item \textbf{Sass: Guide}. \\
        \urlx{https://sass-lang.com/guide/}
  \item \textbf{Bootstrap: Getting Started}. \\
        \urlx{https://getbootstrap.com/docs/5.3/getting-started/introduction/}
  \item \textbf{Get started with Tailwind CSS}. \\
        \urlx{https://tailwindcss.com/docs/installation}
  \item \textbf{Vite: Official Guide}. \\
        \urlx{https://vitejs.dev/guide/}
\end{itemize}
